\section{Spin coefficients and their limits}
\label{sec:spin}
The three-body comparison will use the exact spin spectrum of the wedge map. The four-body comparison will use only finitely many coupling coefficients and their limits.

The spherical pair coefficients can be expressed through angular-momentum
coupling~\cite[Section~II.C and Appendix~A]{OrtizNussinovDukelskySeidel}.
We derive the particular coefficients and fixed-deficit limits needed
here to keep their phases and normalization explicit.

We now fix $e_p$ and $v_p$ to be the normalized lowering bases of $U_Q$ and $V_Q$, respectively, with $v_0=-e_0\wedge e_1$. For spin $j$, lowering from label $p$ to $p+1$ has coefficient $\sqrt{(p+1)(2j-p)}$. Write $(s)_p=s(s-1)\cdots(s-p+1)$, with $(s)_0=1$.

\begin{lemma}[Pair coefficients]\label{lem:pair}
The coefficient of $e_i\wedge e_j$ in $v_p$, for $i<j$, vanishes unless $i+j=p+1$, and otherwise equals
\begin{equation}\label{eq:pair}
 v_p(i,j)=(i-j)\sqrt{\frac{(Q)_i(Q)_j\,p!}{Q(2Q-2)_p\,i!j!}}.
\end{equation}
For fixed $p,i,j$, its limit is
\begin{equation}\label{eq:pairstar}
 v_p^*(i,j)=(i-j)\sqrt{\frac{p!}{2^p i!j!}}\,\mathbf1_{i+j=p+1}.
\end{equation}
\end{lemma}
\begin{proof}
Apply simultaneous lowering $p$ times to $-e_0\wedge e_1$ and divide by $\sqrt{(2Q-2)_p p!}$. Expand the lowering operator binomially. Combining the two contributions to each increasing pair gives \eqref{eq:pair}. Taking the limit gives \eqref{eq:pairstar}.
The normalization can also be checked directly: the square of the
coefficient in \eqref{eq:pair} is
\[
 (i-j)^2\frac{\binom Qi\binom Qj}{Q\binom{2Q-2}{p}},
\]
and differentiating \((1+x)^Q\) gives
\[
 \sum_{i,j=0}^Q(i-j)^2\binom Qi\binom Qj x^{i+j}
 =2Qx(1+x)^{2Q-2}.
\]
Take the coefficient of \(x^{p+1}\) and halve the sum to restrict to \(i<j\).
\end{proof}

The elementary tensor product rule decomposes spins $j_1,j_2$ into one copy of each spin $j_1+j_2-z$, for $0\le z\le\min(2j_1,2j_2)$. This rule can be seen directly in the
polynomial model: at deficit $z$ the raising equation determines a highest
vector, and normalized lowering generates its $2(j_1+j_2-z)+1$
orthogonal weight vectors. Different highest spins give orthogonal
invariant subspaces, and the sum of these dimensions is
$(2j_1+1)(2j_2+1)$, exhausting the tensor product. Let $R_{3,z}$ be the corresponding projection in $V_Q\otimes U_Q$, and put
\[
 d_a=2Q-1,\qquad d_{3,z}=3Q-1-2z.
\]
For $Q\ge2$, the allowed indices are $0\le z\le Q$.

\begin{lemma}[Three-body Gram eigenvalues]\label{lem:q}
\begin{equation}\label{eq:q}
 W_3^\dagger W_3-I=\sum_{z=0}^Qq_z(Q)R_{3,z},\qquad
 q_z(Q)=(-1)^z\frac{(Q)_z}{(2Q-2)_z}
          \left(3z-1-\frac{z(z+1)}Q\right).
\end{equation}
In particular $W_3R_{3,z}=0$ for $z=0,1,3$ whenever these indices are allowed. For fixed $z$,
\begin{equation}\label{eq:qlimit}
 q_z(Q)\longrightarrow (3z-1)(-1/2)^z.
\end{equation}
\end{lemma}
\begin{proof}
Identify $V_Q\otimes U_Q$ with tensors antisymmetric in their first two positions. Then $W_3$ is $\sqrt3$ times full antisymmetrization. Its Gram operator is the compression of $I-P_{13}-P_{23}$, where $P_{ij}$ swaps tensor positions. The two swap compressions coincide, since conjugation by $P_{12}$ interchanges them and $P_{12}=-I$ on the domain. By equivariance and multiplicity-freeness they act by the same scalar $b_z$ on each spin block.

A highest vector of that block has coefficients $s_p$ on $v_p\otimes e_{z-p}$, $0\le p\le z$, satisfying
\begin{equation}\label{eq:highest-ratio}
 \frac{s_{p+1}}{s_p}=-\sqrt{\frac{(z-p)(Q-z+p+1)}{(p+1)(2Q-2-p)}}.
\end{equation}
This follows by applying the raising operator. Write
$\xi_z=\sum_{p=0}^z s_pv_p\otimes e_{z-p}$. Since
$v_0=(-e_0\otimes e_1+e_1\otimes e_0)/\sqrt2$, testing $P_{23}\xi_z$
against $v_0\otimes e_z$ tests $\xi_z$ only on the ordered triples
$(0,z,1)$ and $(1,z,0)$. Their pair labels are $z-1$ and $z$,
respectively, with the first contribution absent at $z=0$.
Substituting their pair coefficients from \eqref{eq:pair} gives
\[
 s_0b_z=\frac12\left[
 s_{z-1}\sqrt{\frac{z(Q)_z}{Q(2Q-2)_{z-1}}}
 -s_z(z-1)\sqrt{\frac{(Q)_z}{(2Q-2)_z}}\right].
\]
Here the ordered tensor coefficients of a normalized wedge include a factor $1/\sqrt2$. Iteration of \eqref{eq:highest-ratio} gives
\[
 \frac{s_z}{s_0}=(-1)^z\sqrt{\frac{(Q)_z}{(2Q-2)_z}},\qquad
 \frac{s_{z-1}}{s_0}=(-1)^{z-1}
 \sqrt{\frac{z(Q-1)_{z-1}}{(2Q-2)_{z-1}}}\quad(z\ge1).
\]
Substitution gives $q_z=-2b_z$ in the form \eqref{eq:q}. The three stated null sectors have $q_z=-1$, and \eqref{eq:qlimit} follows directly.
\end{proof}

We will repeatedly use a fixed-deficit limit of this same calculation. In a product of spins $j_1,j_2$, total lowering deficit means $T=j_1+j_2-m$, where $m$ is the total magnetic quantum number. Coefficients below are with respect to normalized monomials $X^pY^{T-p}/\sqrt{p!(T-p)!}$, and are zero outside the stated index ranges.

\begin{lemma}[Fixed-deficit coupling limit]\label{lem:coupling}
Suppose $j_1,j_2\to\infty$ and $j_2/(j_1+j_2)\to u^2$, where $0<u<1$. Put $v=\sqrt{1-u^2}$. For fixed integers $0\le z\le T$, the normalized coupled vector of spin $j_1+j_2-z$ at total lowering deficit $T$ can be chosen with real coefficients converging to the coefficients $s^{(u)}_{z,T,p}$ of
\begin{equation}\label{eq:coupling}
 \frac{(uX-vY)^z(vX+uY)^{T-z}}{\sqrt{z!(T-z)!}}.
\end{equation}
Descendants in every copy are obtained by normalized lowering, so the matching of bases between equal-spin copies is intertwining.
\end{lemma}
\begin{proof}
At $T=z$, the raising equation gives the adjacent ratio
\[
 -\sqrt{\frac{(z-p)(2j_2-z+p+1)}{(p+1)(2j_1-p)}}.
\]
With sign $(-1)^z$ at $p=0$, normalization yields the coefficients of $(uX-vY)^z/\sqrt{z!}$ in the limit. The change of variables $(X,Y)\mapsto(uX-vY,vX+uY)$ is orthogonal, so these polynomials have norm one in the normalized-monomial inner product. Normalized lowering from $T$ to $T+1$ divides the sum of the two lowering operators by
\[
 \sqrt{(T-z+1)\{2(j_1+j_2-z)-(T-z)\}}.
\]
In the limit this acts by multiplication by $(vX+uY)/\sqrt{T-z+1}$. Induction proves the statement for fixed descendants.
\end{proof}
